
%%FILE blueprint/src/chapter/ap.tex
\chapter{Almost-Periodicity}
\label{chap:ap}


\begin{lemma}[Marcinkiewicz-Zygmund inequality]
\label{mzi}
\lean{Real.marcinkiewicz_zygmund', Real.marcinkiewicz_zygmund}
\leanok
Let $m\geq 1$. If $f:G\to\bbr$ is such that $\bbe_x f(x)=0$ and $\abs{f(x)}\leq 2$ for all $x$ then
\[\bbe_{x_1,\ldots,x_n} \abs{\sum_{i=1}^n f(x_i)}^{2m} \leq (4mn)^{m}.\]
\end{lemma}

\begin{proof}
\leanok
Let $S$ be the left-hand side. Since $0=\bbe_y f(y)$ we have, by the triangle inequality, and H\"{o}lder's inequality,
\begin{align}
  S
  & = \bbe_{x_1, \dots, x_n}\left\lvert \sum_i f(x_i) - \bbe_{y_i} f(y_i)\right\rvert^{2m} \\
  & = \bbe_{x_1, \dots, x_n} \left\lvert\bbe_{y_i}\brac{\sum_i f(x_i) - f(y_i)}\right\rvert^{2m} \\
  & \le \bbe_{x_1,\dots,y_n} \left\lvert \sum_i f(x_i) - f(y_i)\right\rvert^{2m}
\end{align}

Changing the role of $x_i$ and $y_i$ makes no difference here, but multiplies the $i$ summand by $\{-1,+1\}$, and therefore for any $\epsilon_i\in\{-1,+1\}$,
\[ S \leq\bbe_{x_1,\ldots,y_n}\left\lvert \sum_i \epsilon_i(f(x_i)-f(y_i))\right\rvert^{2m}.\]
In particular, if we sample $\epsilon_i\in\{-1,+1\}$ uniformly at random, then
\[ S \leq\bbe_{\epsilon_i} \bbe_{x_1,\ldots,y_n}\left\lvert \sum_i \epsilon_i(f(x_i)-f(y_i))\right\rvert^{2m}.\]
We now change the order of expectation and consider the expectation over just $\epsilon_i$, viewing the $f(x_i)-f(y_i)=z_i$, say, as fixed quantities. For any $z_i$ we can expand $\bbe_{\epsilon_i} \lvert \sum_i \epsilon_iz_i\rvert^{2m}$ and then bound it from above, using the triangle inequality and $\abs{z_i}\leq 4$, by
\[4^{2m}\sum_{k_1+\cdots+k_n=2m}\binom{2m}{k_1,\ldots,k_n}\abs{\bbe\epsilon_1^{k_1}\cdots \epsilon_n^{k_n}}.\]
The inner expectation vanishes unless each $k_i$ is even, when it is trivially one. Therefore the above quantity is exactly
\[\sum_{l_1+\cdots+l_n=m}\binom{2m}{2l_1,\ldots,2l_n}\leq m^mn^m,\]
since for any $l_1+\cdots+l_n=m$,
\[\binom{2m}{2l_1,\ldots,2l_n}\leq m^m\binom{m}{l_1,\ldots,l_n}.\]
This can be seen, for example, by writing both sides out using factorials, yielding

\[\frac{(2m)!}{(2l_1)!\cdots (2l_n)!}\leq \frac{(2m)!}{2^mm!}\frac{m!}{l_1!\cdots l_n!}\leq m^m\frac{m!}{l_1!\cdots l_n!}.\]
\end{proof}


\begin{lemma}[Complex-valued Marcinkiewicz-Zygmund inequality]
\label{mzi_complex}
\leanok
\lean{RCLike.marcinkiewicz_zygmund}
Let $m\geq 1$. If $f:G\to\bbc$ is such that $\bbe_x f(x)=0$ and $\abs{f(x)}\leq 2$ for all $x$ then
\[\bbe_{x_1,\ldots,x_n} \abs{\sum_{i=1}^n f(x_i)}^{2m} \leq (8mn)^{m}.\]
\end{lemma}

\begin{proof}
\uses{mzi}
\leanok
Test.
\end{proof}


\begin{lemma}
  \label{random_approx_expect}
  \lean{AlmostPeriodicity.lemma28}
  \leanok

  Let $\epsilon>0$ and $m\geq 1$. Let $A\subseteq G$ and $f:G\to \bbc$.
  If $k\geq 64m\epsilon^{-2}$ then the set
  \[
    L = \{\tup{a}\in A^k : \|\tfrac{1}{k}\sum_{i=1}^k f(x-a_i)-\mu_A\ast f\|_{2m}
      \le \epsilon \| f\|_{2m}\}.
  \]
  has size at least $\lvert A \rvert^k/2$.
\end{lemma}

\begin{proof}
\uses{mzi_complex}
\leanok
Note that if $a\in A$ is chosen uniformly at random then, for any fixed $x\in G$,
\[\bbe f(x-a_i)= \frac{1}{\abs{A}}\sum_{a\in A}f(x-a)=\frac{1}{\abs{A}}\ind{A}\ast f(x)=\mu_A\ast f(x).\]
Therefore, if we choose $a_1,\ldots,a_k\in A$ independently uniformly at random, for any fixed $x\in G$ and $1\leq i\leq k$,  the random variable $f(x-a_i)-f\ast \mu_A(x)$ has mean zero. By the Marcinkiewicz-Zygmund inequality Lemma~\ref{mzi}, therefore,
\begin{multline*}
\bbe\abs{ \frac{1}{k}\sum_i f(x-a_i)-f\ast \mu_A(x)}^{2m} \leq \\(16m/k)^mk^{-1} \bbe \sum_i \abs{f(x-a_i)-f\ast \mu_A(x)}^{2m}.
\end{multline*}
We now sum both sides over all $x\in G$. By the triangle inequality, for any fixed $1\leq i\leq k$ and $a_i\in A$,
\begin{align*}
\sum_{x\in G} \abs{f(x-a_i)-f\ast \mu_A(x)}^{2m}
&\leq 2^{2m-1}\sum_{x\in G}\abs{f(x-a_i)}^{2m}+\sum_{x\in G}\abs{f\ast \mu_A(x)}^{2m}\\
&\leq 2^{2m-1}\brac{\norm{f}_{2m}^{2m}+\norm{f\ast \mu_A}_{2m}^{2m}}.
\end{align*}
We note that $\norm{\mu_A}_1=\frac{1}{\abs{A}}\sum_{x\in A}\ind{A}(x)=\abs{A}/\abs{A}=1$, and hence by Young's inequality, $\norm{f\ast \mu_A}_{2m}\leq \norm{f}_{2m}$, and so
\[\sum_{x\in G} \abs{f(x-a_i)-f\ast \mu_A(x)}^{2m}\leq 2^{2m}\norm{f}_{2m}^{2m}.\]
It follows that
\[\bbe_{a_1,\ldots,a_k\in A}\norm{\frac{1}{k}\sum_i\tau_{a_i}f-f\ast \mu_A}_{2m}^{2m}\leq
(64m/k)^m\norm{f}_{2m}^{2m}.\]
In particular, if $k\geq 64\epsilon^{-2}m$ then the right-hand side is at most $(\frac{\epsilon}{2}\norm{f}_{2m})^{2m}$ as required.
\end{proof}


\begin{lemma}
  \label{aps_in_translates}
  \lean{AlmostPeriodicity.just_the_triangle_inequality}
  \leanok

  Let $A\subseteq G$ and $f:G\to \bbc$. Let $\epsilon>0$ and $m\geq 1$ and $k\geq 1$. Let
  \[
    L = \{\tup{a}\in A^k : \|\tfrac{1}{k}\sum_{i=1}^k f(x-a_i) - \mu_A\ast f\|_{2m}
      \leq \epsilon \| f\|_{2m}\}.
  \]
  If $t\in G$ is such that $\tup{a}\in L$ and $\tup{a}+(t,\ldots,t)\in L$ then
  \[\| \tau_t(\mu_A\ast f)-\mu_A\ast f\|_{2m}\leq 2\epsilon \|f\|_{2m}.\]
\end{lemma}

\begin{proof}
\leanok
Test.
\end{proof}


\begin{lemma}
\label{lots_of_diagonals}
\lean{big_shifts}
\leanok
Let $A\subseteq G$ and $k\geq 1$ and $L\subseteq A^k$. Then there exists some $\tup{a}\in L$ such that
\[\#\{ t\in G : \tup{a}+(t,\ldots,t)\in L\}\geq \frac{\lvert L\rvert}{\lvert A+S\rvert^k}\lvert S\rvert.\]
\end{lemma}

\begin{proof}
\leanok
Test.
\end{proof}


\begin{theorem}[$L_p$ almost periodicity]
\label{lp_ap}
\lean{AlmostPeriodicity.almost_periodicity}
\leanok
Let $\epsilon\in (0,1]$ and $m\geq 1$. Let $K\geq 2$ and $A,S\subseteq G$ with $\lvert A+S\rvert\leq K\lvert A\rvert$.
Let $f:G\to \bbc$. There exists $T\subseteq G$ such that
\[\lvert T\rvert \geq K^{-512m\epsilon^{-2}}\lvert S\rvert\]
such that for any $t\in T$ we have
\[\| \tau_t(\mu_A\ast f)-\mu_A\ast f\|_{2m}\leq \epsilon \| f\|_{2m}.\]
\end{theorem}

\begin{proof}
\uses{random_approx_expect, lots_of_diagonals, aps_in_translates}
\leanok
Test.
\end{proof}


\begin{theorem}[$L_\infty$ almost periodicity]
\label{linfty_ap}
\lean{AlmostPeriodicity.linfty_almost_periodicity}
\leanok
Let $\epsilon\in (0,1]$. Let $K\geq 2$ and $A,S\subseteq G$ with $\lvert A+S\rvert\leq K\lvert A\rvert$.
Let $B,C\subseteq G$. Let $\eta=\min(1,\lvert C\rvert/\lvert B\rvert)$. There exists $T\subseteq G$ such that
\[\lvert T\rvert \geq K^{-4096\lceil \lo{\eta}\rceil\epsilon^{-2}}\lvert S\rvert\]
such that for any $t\in T$ we have
\[\| \tau_t(\mu_A\ast 1_B\ast \mu_C)-\mu_A\ast 1_B\ast \mu_C\|_{\infty}\leq \epsilon.\]
\end{theorem}

\begin{proof}
\leanok
\uses{lp_ap} Let $T$ be as given in \ref{lp_ap}
with $f=1_B$ and $m=\lceil \lo{\eta}\rceil$ and $\epsilon=\epsilon/e$. (The size bound on $T$ follows since $e^2\leq 8$.) Fix $t\in T$ and let $F=\tau_t(\mu_A\ast 1_B)-\mu_A\ast 1_B$. We have, for any $x\in G$,
\[(\tau_t(\mu_A\ast 1_B\ast \mu_C)-\mu_A\ast 1_B\ast \mu_C)(x)=F\ast \mu_C(x)=\sum_y F(y)\mu_{C}(x-y)=\sum_yF(y)\mu_{x-C}(y).\]
By Hölder's inequality, this is (in absolute value), for any $m\geq 1$,
\[\norm{F}_{2m}\norm{\mu_{x-C}}_{1+\frac{1}{2m-1}}.\]
By the construction of $T$ the first factor is at most
$\frac{\epsilon}{e}\| 1_B\|_{2m}=\frac{\epsilon}{e}\lvert B\rvert^{1/2m}$.
We have by calculation
\[\norm{\mu_{x-C}}_{1+\frac{1}{2m-1}}=\lvert x-C\rvert^{-1/2m}=\lvert C\rvert^{-1/2m}.\]
Therefore we have shown that

\[\| \tau_t(\mu_A\ast 1_B\ast \mu_C)-\mu_A\ast 1_B\ast \mu_C\|_{\infty}\leq \frac{\epsilon}{e}(\lvert C\rvert/\lvert B\rvert)^{-1/2m}.\]
The claim now follows since, by choice of $m$,
\[(\lvert C\rvert/\lvert B\rvert)^{-1/2m}\leq e\]
(dividing into cases as to whether $\eta=1$ or not).
\end{proof}


\begin{theorem}
\label{linfty_ap_boosted}
\lean{AlmostPeriodicity.linfty_almost_periodicity_boosted}
\leanok
Let $\epsilon\in (0,1]$ and $k\geq 1$. Let $K\geq 2$ and $A,S\subseteq G$ with $\lvert A+S\rvert\leq K\lvert A\rvert$.
Let $B,C\subseteq G$. Let $\eta=\min(1,\lvert C\rvert/\lvert B\rvert)$. There exists $T\subseteq G$ such that
\[\lvert T\rvert \geq K^{-4096\lceil \lo{\eta}\rceil k^2\epsilon^{-2}}\lvert S\rvert\]
such that
\[\| \mu_T^{(k)}\ast \mu_A\ast 1_B\ast \mu_C-\mu_A\ast 1_B\ast \mu_C\|_{\infty}\leq \epsilon.\]
\end{theorem}

\begin{proof}
\uses{linfty_ap}
\leanok
Let $T$ be as in Theorem~\ref{linfty_ap} with $\epsilon$ replaced by $\epsilon/k$. Note that, for any $x\in G$,
\[\mu_T^{(k)}\ast \mu_A\ast 1_B\ast \mu_C(x)=\frac{1}{\lvert T\rvert^k}\sum_{t_1,\ldots,t_k\in T}\tau_{t_1+\cdots+t_k}\mu_A\ast 1_B\ast \mu_C(x).\]
It therefore suffices (by the triangle inequality) to show, for any fixed $x\in G$ and $t_1,\ldots,t_k\in T$, that with $F=\mu_A\ast 1_B\ast \mu_C$, we have
\[\lvert \tau_{t_1+\cdots+t_k}F(x)-F(x)\rvert \leq \epsilon.\]
This follows by the triangle inequality applied $k$ times if we knew that, for $1\leq l\leq k$,
\[\lvert \tau_{t_1+\cdots+t_l}F(x)-\tau_{t_1+\cdots+t_{l-1}}F(x)\rvert \leq \epsilon/k.\]
We can write the left-hand side as
\[\lvert \tau_{t_1+\cdots+t_l}F(x)-\tau_{t_1+\cdots+t_{l-1}}F(x)\rvert=\lvert \tau_{t_l}F(x-t_1-\cdots-t-{l-1})-F(x-t_1-\cdots-t-{l-1})\rvert.\]
The right-hand side is at most
\[\| \tau_{t_l}F-F\|_\infty\]
and we are done by choice of $T$.
\end{proof}


%%FILE blueprint/src/chapter/bohr.tex
\chapter{Bohr sets}

\begin{definition}[Bohr sets]
\label{bohr-set}
\lean{BohrSet}
\leanok
Let $\nu: \widehat{G} \to \mathbb{R}$. The corresponding Bohr set is defined to be
\[\mathrm{Bohr}(\nu)=\left\{ x\in G : \abs{1-\gamma(x)}\leq \nu(\gamma)\textrm{ for all }\gamma\in\Gamma\right\}.\]
The rank of $\nu$, denoted by $\mathrm{rk}(\nu)$, is defined to be the size of the set of those $\gamma\in\widehat{G}$ such that $\nu(\gamma)<2$.

(Basic API facts: Bohr sets are symmetric and contain $0$. Also that, without loss of generality, we can assume $\nu$ takes only values in $\mathbb{R}_{\geq 0}$ - I think it might be easier to have the definition allow arbitrary real values, and then switch to non-negative only in proofs where convenient. Or could have the definition only allow non-negative valued functions in the first place.)
\end{definition}

\begin{lemma}
  \label{bohr-size}
  \uses{bohr-set}
  % \lean{BohrSet.le_card_smul}
  % \leanok

  If $\rho\in (0,1)$ and $\nu:\widehat{G}\to \mathbb{R}$ then
  \[\abs{\mathrm{Bohr}(\rho\cdot \nu)}\geq (\rho/4)^{\mathrm{rk}(\nu)}\abs{\mathrm{Bohr}(\nu)}.\]
\end{lemma}
\begin{proof}
  There are at most $\lceil 4/\rho\rceil$ many $z_i$ such that if $\lvert 1-w\rvert \leq \nu(\gamma)$ then $\lvert z_i-w\rvert\leq \rho\nu(\gamma)/2$ for some $i$.
  Let $\Gamma=\{ \gamma : \nu(\gamma) <2\}$
  and define a function $f:\mathrm{Bohr}(\nu)\to \lceil 2/\rho\rceil^{\mathrm{rk}(\nu)}$
  where for $\gamma\in \Gamma$ we assign the $\gamma$-coordinate of $f(x)$
  as whichever $j$ has $\lvert z_j-\gamma(x)\rvert \leq \rho\nu(\gamma)/2$.

  By the pigeonhole principle there must exist some $(j_1,\ldots,j_{d})$ such that
  $f^{-1}(j_1,\ldots,j_d)$ has size at least
  $(\lceil 2/\rho\rceil)^{-\mathrm{rk}(\nu)}\abs{\mathrm{Bohr}(\nu)}$. Call this set $B'$.
  It must be non-empty, so fix some $x\in B'$.
  We claim that $B'-x\subseteq \abs{\mathrm{Bohr}(\rho\cdot \nu)}$, which completes the proof.

  Suppose that $z=x+y$ with $x,y\in B'$, and fix some $\gamma\in \Gamma$.
  By assumption there is some $z_j\in \mathbb{C}$ such that
  $\lvert z_j-\gamma(x)\rvert \leq \rho \nu(\gamma)/2$ and
  $\lvert z_j-\gamma(y)\rvert \leq \rho \nu(\gamma)/2$. Then by the triangle inequality,
  \[\lvert 1-\gamma(y-x)\rvert=\lvert \gamma(x)-\gamma(y)\rvert \leq \rho \nu(\gamma)\]
  and so $z=y-x\in \mathrm{Bohr}(\rho\cdot \nu)$.
\end{proof}

\begin{definition}[Regularity]
\label{bohr-reg-def}
\uses{bohr-set}
\lean{BohrSet.IsRegular}
\leanok
We say $\nu:\widehat{G}\to \mathbb{R}$ is regular if, with $d=\mathrm{rk}(\nu)$, for all $\kappa\in\mathbb{R}$ with $\abs{\kappa}\leq 1/100d$ we have
\[(1-100 d\abs{\kappa})\leq \frac{\abs{\mathrm{Bohr}((1+\kappa)\nu)}}{\abs{\mathrm{Bohr}(\nu)}}\leq(1+ 100 d\abs{\kappa})\]
\end{definition}

\begin{lemma}
\label{bohr-regularity}
\uses{bohr-reg-def}
\lean{BohrSet.regularity}
\leanok
For any $\nu:\widehat{G}\to \mathbb{R}$ there exists $\rho\in[\frac{1}{2},1]$ such that $\rho\cdot \nu$ is regular.
\end{lemma}
\begin{proof}
\uses{bohr-size}
To do.
\end{proof}

\begin{lemma}
\label{reg-conv}
\uses{bohr-reg-def}
If $B$ is a regular Bohr set of rank $d$ and $\mu:G\to\bbr_{\geq 0}$ is supported on $B_\rho$, with $\rho \in (0,1)$, then
\[ \norm{ \mu_B*\mu - \mu_B }_{1} \ll \rho d\norm{\mu}_1. \]
\end{lemma}
\begin{proof}
To do.
\end{proof}

\begin{lemma}
\label{bohr-majorise}
\uses{bohr-reg-def}
There is a constant $c>0$ such that the following holds.  Let $B$ be a regular Bohr set of rank $d$ and $L\geq 1$ be any integer. If $\nu:G\to\bbr_{\geq 0}$ is supported on $L B_\rho$, where $\rho \leq c/Ld$, and $\norm{\nu}_1=1$, then
\[\mu_B \leq 2\mu_{B_{1+L\rho}}\ast \nu.\]
\end{lemma}
\begin{proof}
\uses{reg-conv}
To do.
\end{proof}

\begin{lemma}
\label{bourgain-trick}
\uses{bohr-reg-def}
There is a constant $c>0$ such that the following holds. Let $B$ be a regular Bohr set of rank $d$, suppose $A\subseteq B$ has density $\alpha$, let $\epsilon>0$, and suppose $B',B''\subseteq B_\rho$ where $\rho\leq c\alpha\epsilon/d$. Then either
\begin{enumerate}
\item there is some translate $A'$ of $A$ such that $\abs{A'\cap B'}\geq (1-\epsilon)\alpha\abs{B'}$ and $\abs{A'\cap B''}\geq (1-\epsilon)\alpha \abs{B''}$, or
\item $\norm{\ind{A}\ast \mu_{B'}}_\infty\geq (1+\epsilon/2)\alpha$, or
\item $\norm{\ind{A}\ast \mu_{B''}}_\infty \geq (1+\epsilon/2)\alpha$.
\end{enumerate}
\end{lemma}
\begin{proof}
\uses{reg-conv}
To do.
\end{proof}


%%FILE blueprint/src/chapter/chang.tex
\chapter{Chang's lemma}

\begin{definition}[Dissociation]
\label{dissociated}
\lean{AddDissociated}
\leanok
We say that $A\subseteq G$ is dissociated if, for any $m\geq 1$, and any $x\in G$, there is at most one $A'\subset A$ of size $\abs{A'}=m$ such that
\[\sum_{a\in A'}a=x.\]
\end{definition}


\begin{lemma}[Rudin's exponential inequality]
\label{rudin_exp}
\uses{dissociated}
\lean{rudin_exp_ineq}
\leanok
If the discrete Fourier transform of $f : G \longrightarrow \C$ has dissociated support, then
$$\E \exp(\Re f) \le \exp\left(\frac{\norm f_2^2} 2\right)$$
It follows that
\[\E_x e^{\abs{f(x)}} \le 2e^{\| f\|_2^2/2}.\]
\end{lemma}
\begin{proof}
\uses{mzi_complex}
\leanok
Using the convexity of $t\mapsto e^{tx}$ (for all $x\geq 0$ and $t\in[-1,1]$) we have
\[e^{tx}\le \cosh(x)+t\sinh(x).\]
It follows (taking $x=\lvert z\rvert$ and $t=\Re(z)/\lvert z\rvert$) that, for any $z\in \mathbb{C}$,
\[e^{\Re z}\le \cosh \lvert z\rvert +\Re(z/\lvert z\rvert)\sinh \lvert z\rvert.\]
In particular, if $c_\gamma\in \mathbb{C}$ with $\lvert c_\gamma\rvert=1$ is such that $\widehat{f}(\gamma)=c_\gamma\lvert \widehat{f}(\gamma)\rvert$, then
\begin{align*}
e^{\Re f(x)}
&= \exp\left( \Re \sum_{\gamma\in\Gamma}\widehat{f}(\gamma)\gamma(x)\right)\\
&=\prod_{\gamma\in \Gamma} \exp\left( \Re \widehat{f}(\gamma)\gamma(x)\right)\\
&\le \prod_{\gamma\in \Gamma}\left( \cosh \lvert \widehat{f}(\gamma)\rvert +\Re c_\gamma \gamma(x)\sinh \lvert \widehat{f}(\gamma)\rvert\right).
\end{align*}
Therefore
\[\E_x e^{\Re f(x)}\le \E_x \prod_{\gamma\in \Gamma} \left( \cosh\lvert \widehat{f}(\gamma)\rvert+\Re c_\gamma \gamma(x)\sinh\lvert \widehat{f}(\gamma)\rvert\right).\]
Using $\Re z=(z+\overline{z})/2$ the product here can be expanded as the sum of
\[\prod_{\gamma\in \Gamma_2}\frac{c_\gamma}{2}\prod_{\gamma\in \Gamma_3}\frac{\overline{c_\gamma}}{2}\left(\prod_{\gamma\in \Gamma_1}\cosh\lvert \widehat{f}(\gamma)\rvert\right)\left(\prod_{\gamma\in \Gamma_2\cup\Gamma_3}\sinh\lvert \widehat{f}(\gamma)\rvert\right)\left(\sum_{\gamma\in \Gamma_2}\gamma-\sum_{\lambda\in \Gamma_3}\lambda\right)(x)\]
as $\Gamma_1\sqcup \Gamma_2\sqcup \Gamma_3=\Gamma$ ranges over all partitions of $\Gamma$ into three disjoint parts. Using the definition of dissociativity we see that
\[\sum_{\gamma\in \Gamma_2}\gamma-\sum_{\lambda\in \Gamma_3}\lambda\neq 0\]
unless $\Gamma_2=\Gamma_3=\emptyset$. In particular summing this term over all $x\in G$ gives $0$. Therefore the only term that survives averaging over $x$ is when $\Gamma_1=\Gamma$, and so
\[\E_x e^{\Re f(x)}\le \prod_{\gamma\in \Gamma} \cosh \lvert \widehat{f}(\gamma)\rvert.\]
The conclusion now follows using $\cosh(x) \le e^{x^2/2}$ and $\sum_{\gamma\in \Gamma}\lvert \widehat{f}(\gamma)\rvert^2=\| f\|_2^2$. The second conclusion follows by applying it to $f(x)$ and $-f(x)$ and using
\[e^{\abs{y}}\le e^y+e^{-y}.\]
\end{proof}


\begin{lemma}[Rudin's inequality]
\label{rudin}
\uses{dissociated}
\lean{rudin_ineq}
\leanok
If the discrete Fourier transform of $f : G \longrightarrow \C$ has dissociated support and $p \ge 2$ is an integer, then $\norm{f}_p \le 4\sqrt{pe} \norm f_2$.
\end{lemma}
\begin{proof}
\uses{rudin_exp}
\leanok
It is enough to show that $\norm{\Re f}_p \le 2\sqrt{pe} \norm f_2$ as then
$$\norm{f}_p \le \norm{\Re f}_p + \norm{i \Im f}_p = \norm{\Re f}_p + \norm{\Re (-if)}_p \le 4\sqrt{pe} \norm f_2$$

If $f = 0$, the result is obvious. So assume $f \ne 0$. $\norm{f}_2 > 0$, so WLOG $\norm{f}_2 = \sqrt p$.

Rudin's exponential inequality for $f$ becomes $\E \exp|\Re f| \le 2\exp(\frac p 2) = (2\sqrt e)^p$. Using $\frac{x^p}{p!} \le e^x$ for positive $x$, we get
$$\frac{\norm{\Re f}_p^p}{p^p} \le \frac{\norm{\Re f}_p^p}{p!} = \E \frac{|\Re f|^p}{p!} \le \E \exp|\Re f|$$
Rearranging, $\norm{\Re f}_p \le 2 p\sqrt{e} = 2 \sqrt{pe} \norm f_2$.
% TODO: Normalisation issue
\end{proof}


\begin{definition}[Large spectrum]
  \label{large_spec}
  \lean{largeSpec}
  \leanok

  Let $G$ be a finite abelian group and $f:G\to\bbc$. Let $\eta\in \bbr$. The $\eta$-large spectrum is defined to be
  \[\Delta_\eta(f) = \{ \gamma\in\widehat{G} : \lvert \widehat{f}(\gamma)\rvert \geq \eta\norm{f}_1\}.\]
\end{definition}


\begin{definition}[Weighted energy]
\label{weight_energy}
\lean{energy}
\leanok
Let $\Delta\subseteq \widehat{G}$ and $m\geq 1$. Let $\nu:G\to \bbc$. Then
\[E_{2m}(\Delta;\nu)=\sum_{\gamma_1,\ldots,\gamma_{2m}\in \Delta}\abs{\widehat{\nu}(\gamma_1+\cdots-\gamma_{2m})}.\]
\end{definition}


\begin{definition}[Energy]
\label{energy}
\uses{weight_energy}
\lean{boringEnergy}
\leanok
Let $G$ be a finite abelian group and $A\subseteq G$. Let $m\geq 1$. We define
\[E_{2m}(A)=\sum_{a_1,\ldots,a_{2m}\in A}1_{a_1+\cdots-a_{2m}=0}.\]
\end{definition}


\begin{lemma}
\label{general_hoelder}
\uses{large_spec, weight_energy}
\lean{general_hoelder}
\leanok
Let $G$ be a finite abelian group and $f:G\to\bbc$. Let $\nu:G\to \bbr_{\geq 0}$ be such that whenever $\abs{f}\neq 0$ we have $\nu \geq 1$. Let $\Delta\subseteq \Delta_\eta(f)$. Then, for any $m\geq 1$.
\[\eta^{2m}\frac{\norm{f}_1^2}{\norm{f}_2^2}\abs{\Delta}^{2m}\le E_{2m}(\Delta;\nu).\]
\end{lemma}

\begin{proof}
\leanok
By definition of $\Delta_\eta(f)$ we know that
\[\eta\norm{f}_1\abs{\Delta}\le \sum_{\gamma\in\Delta} \lvert \widehat{f}(\gamma)\rvert.\]
There exists some $c_\gamma\in\bbc$ with $\lvert c_\gamma\rvert=1$ for all $\gamma$ such that
\[\lvert \widehat{f}(\gamma)\rvert=c_\gamma\widehat{f}(\gamma)=c_\gamma \sum_{x\in G}f(x)\overline{\gamma(x)}.\]
Interchanging the sums, therefore,

\[\eta\norm{f}_1\abs{\Delta}\le \sum_{x\in G}f(x)\sum_{\gamma\in\Delta} c_\gamma \overline{\gamma(x)}.\]
By H\"{o}lder's inequality the right-hand side is at most
\[\brac{\sum_x \abs{f(x)}}^{1-1/m}\brac{\sum_x \abs{f(x)}\abs{\sum_{\gamma\in\Delta}c_\gamma\overline{\gamma(x)}}^m}^{1/m}.\]
Taking $m$th powers, therefore, we have
\[\eta^m\abs{\Delta}^m\norm{f}_1\le \sum_{x}\abs{f(x)}\abs{\sum_{\gamma\in\Delta}c_\gamma\overline{\gamma(x)}}^m.\]
By assumption we can bound $\abs{f(x)}\le \abs{f(x)}\nu(x)^{1/2}$, and therefore by the Cauchy-Schwarz inequality the right-hand side is bounded above by
\[\norm{f}_2\brac{\sum_x \nu(x)\abs{\sum_{\gamma\in\Delta}c_\gamma\overline{\gamma(x)}}^{2m}}^{1/2}.\]
Squaring and simplifying, we deduce that
\[\eta^{2m}\abs{\Delta}^{2m}\frac{\norm{f}_1^2}{\norm{f}_2^2}\le \sum_x \nu(x)\abs{\sum_{\gamma\in\Delta}c_\gamma\overline{\gamma(x)}}^{2m}.\]
Expanding out the power, the right-hand side is equal to
\[\sum_x \nu(x)\sum_{\gamma_1,\ldots,\gamma_{2m}}c_{\gamma_1}\cdots \overline{c_{\gamma_{2m}}} (\overline{\gamma_1}\cdots \gamma_{2m})(x).\]
Changing the order of summation this is equal to
\[\sum_{\gamma_1,\ldots,\gamma_{2m}}c_{\gamma_1}\cdots \overline{c_{\gamma_{2m}}}
\widehat{\nu}(\gamma_1\cdots \overline{\gamma_{2m}}).\]
The result follows by the triangle inequality.
\end{proof}


\begin{lemma}
\label{spec_hoelder}
\uses{energy}
\lean{spec_hoelder}
\leanok
Let $G$ be a finite abelian group and $f:G\to\bbc$. Let $\Delta\subseteq \Delta_\eta(f)$. Then, for any $m\geq 1$.
\[N^{-1}\eta^{2m}\frac{\norm{f}_1^2}{\norm{f}_2^2}\abs{\Delta}^{2m}\le E_{2m}(\Delta).\]
\end{lemma}

\begin{proof}
\uses{general_hoelder}
\leanok
Apply Lemma~\ref{general_hoelder} with $\nu\equiv 1$, and use the fact that $\sum_x \lambda(x)$ is $N$ if $\lambda\equiv 1$ and $0$ otherwise.
\end{proof}


\begin{lemma}
\label{energy_alt}
\uses{energy}
\lean{boringEnergy_eq}
\leanok
If $A\subset G$ and $m\geq 1$ then
\[E_{2m}(A) = \sum_x 1_A^{(m)}(x)^2.\]
\end{lemma}

\begin{proof}
\leanok
Expand out definitions.
\end{proof}


\begin{lemma}
\label{diss_energy}
\uses{dissociated, energy}
\lean{AddDissociated.boringEnergy_le}
\leanok
If $A\subseteq G$ is dissociated then $E_{2m}(A) \le (32e m \abs{A})^m$.
\end{lemma}
\begin{proof}
\uses{energy_alt, rudin}
\leanok
By Lemma~\ref{energy_alt} and Lemma~\ref{rudin}
\begin{align*}
  E_{2m}(A)
  & = \E_\gamma \abs{\hat 1_A(\gamma)}^{2m} \\
  & = \norm{\hat 1_A}_{2m}^{2m} \\
  & \le (4\sqrt{2em})^{2m} \norm{\hat 1_A}_2^{2m} \\
  & = (32em)^m \norm{1_A}_2^{2m} \\
  & = (32em)^m \abs{A}^m
\end{align*}
\end{proof}


\begin{lemma}
  \label{diss_span}
  \uses{dissociated}
  \lean{Finset.exists_subset_addSpan_card_le_of_forall_addDissociated}
  \leanok

  If $A\subseteq G$ contains no dissociated set with $\geq K+1$ elements,
  then there is $A'\subseteq A$ of size $\abs{A'}\le K$ such that
  \[A\subseteq \left\{ \sum_{a\in A'}c_aa : c_a\in \{-1,0,1\} \right\}.\]
\end{lemma}

\begin{proof}
\leanok
Let $A'\subseteq A$ be a maximal dissociated subset (this exists and is non-empty, since trivially any singleton is dissociated). We have $\abs{A'}\le K$ by assumption.

Let $S$ be the span on the right-hand side. It is obvious that $A'\subseteq S$. Suppose that $x\in A\backslash A'$. Then $A'\cup\{x\}$ is not dissociated by maximality. Therefore there exists some $y\in G$ and two distinct sets $B,C\subseteq A'\cup \{x\}$ such that
\[\sum_{b\in B}b = y = \sum_{c\in C} c.\]
If $x\not\in B$ and $x\not\in C$ then this contradicts the dissociativity of $A'$. If $x\in B$ and $x\in C$ then we have
\[\sum_{b\in B\backslash x}b=y-x=\sum_{c\in C\backslash x}c,\]
again contradicting the dissociativity of $A'$. Without loss of generality, therefore, $x\in B$ and $x\not\in C$. Therefore
\[x=\sum_{c\in C}c - \sum_{b\in B\backslash x}b\]
which is in the span as required.
\end{proof}


\begin{theorem}[Chang's lemma]
\label{chang}
\uses{large_spec}
\lean{chang}
\leanok
Let $G$ be a finite abelian group and $f:G\to \bbc$. Let $\eta >0$ and $\alpha=N^{-1}\norm{f}_1^2/\norm{f}_2^2$. There exists some $\Delta\subseteq \Delta_\eta(f)$ such that
\[\abs{\Delta}\le \lceil e\mathcal{L}(\alpha)\eta^{-2}\rceil \]
and
\[\Delta_\eta(f)\subseteq \left\{ \sum_{\gamma\in\Delta}c_\gamma \gamma : c_\gamma\in \{-1,0,1\} \right\}.\]
\end{theorem}

\begin{proof}
\uses{diss_energy, diss_span, spec_hoelder}
\leanok
By Lemma~\ref{diss_span} it suffices to show that $\Delta_\eta(f)$ contains no dissociated set with at least
\[K= \lceil e\mathcal{L}(\alpha)\eta^{-2}\rceil+1\]
many elements. Suppose not, and let $\Delta\subseteq \Delta_\eta(f)$ be a dissociated set of size $K$. Then by Lemma~\ref{diss_energy} we have, for any $m\geq 1$,
\[E_{2m}(\Delta)\le m!K^m.\]
On the other hand, by Lemma~\ref{spec_hoelder},

\[\eta^{2m}\alpha K^{2m}\le E_{2m}(\Delta).\]
Rearranging these bounds, we have
\[K^m \le m! \alpha^{-1}\eta^{-2m}\le m^m\alpha^{-1}\eta^{-2m}.\]
Therefore $K\le \alpha^{-1/m}m\eta^{-2}$. This is a contradiction to the choice of $K$ if we choose $m=\mathcal{L}(\alpha)$, since $\alpha^{-1/m}\le e$.
\end{proof}


%%FILE blueprint/src/chapter/drc.tex
\chapter{Dependent random choice}


\begin{lemma}
\label{drc}
\lean{drc}
\leanok
Let $p\geq 2$ be an even integer. Let $B_1,B_2\subseteq G$ and $\mu=\mu_{B_1}\circ\mu_{B_2}$. For any finite set $A\subseteq G$ and function $f:G\to\bbr_{\geq 0}$ there exist $A_1\subseteq B_1$ and $A_2\subseteq B_2$ such that
\[\langle \mu_{A_1}\circ \mu_{A_2}, f\rangle \norm{\ind{A}\circ \ind{A}}_{p(\mu)}^p\leq 2\langle (\ind{A}\circ \ind{A})^p,f\rangle_\mu\]
and
\[\min\brac{\frac{\abs{A_1}}{\abs{B_1}},\frac{\abs{A_2}}{\abs{B_2}}}\geq \frac{1}{4}\abs{A}^{-2p}\norm{\ind{A}\circ \ind{A}}_{p(\mu)}^{2p}.\]
\end{lemma}

\begin{proof}
\leanok
First note that the statement is trivially true (with $A_1=B_1$ and $A_2=B_2$, say) if $\norm{\ind{A}\circ \ind{A}}_{p(\mu)}^p=0$. We can therefore assume this is $\neq 0$.

For $s\in G^{p}$ let $A_1(s)=B_1\cap (A+s_1)\cap\cdots\cap (A+s_{p})$, and similarly for $A_2(s)$. Note that
\begin{align*}
\langle (\ind{A}\circ \ind{A})^p, f\rangle_\mu
&=\sum_{x} \mu_{B_1}\circ \mu_{B_2}(x)(\ind{A}\circ \ind{A}(x))^pf(x)\\
&=\sum_{b_1,b_2}\mu_{B_1}(b_1)\mu_{B_2}(b_2)\ind{A}\circ \ind{A}(b_1-b_2)^{p}f(b_1-b_2)\\
&=\sum_{b_1,b_2}\mu_{B_1}(b_1)\mu_{B_2}(b_2)\brac{\sum_{t\in G}\ind{A+t}(b_1)\ind{A+t}(b_2)}^{p}f(b_1-b_2)\\
&=\sum_{b_1,b_2}\mu_{B_1}(b_1)\mu_{B_2}(b_2)\sum_{s\in G^p}\ind{A_1(s)}(b_1)\ind{A_2(s)}(b_2)f(b_1-b_2)\\
&=\abs{B_1}^{-1}\abs{B_2}^{-1}\sum_{s\in G^p}\langle \ind{A_1(s)}\circ \ind{A_2(s)}, f\rangle.
\end{align*}
In particular, applying this with $f\equiv 1$ we see that
\[\norm{\ind{A}\circ \ind{A}}_{p(\mu)}^p=\abs{B_1}^{-1}\abs{B_2}^{-1}\sum_s\abs{A_1(s)}\abs{A_2(s)}\]
and
\[\frac{\langle (\ind{A}\circ \ind{A})^p,f\rangle_\mu}{\norm{\ind{A}\circ \ind{A}}_{p(\mu)}^p}=\frac{\sum_{s}\langle \ind{A_1(s)}\circ \ind{A_2(s)}, f\rangle}{\sum_s \abs{A_1(s)}\abs{A_2(s)}}=\eta,\]
say. Let $M>0$ be some parameter, and let
\[g(s) = \begin{cases} 1&\textrm{ if }0<\abs{A_1(s)}\abs{A_2(s)}<M^2\textrm{ and }\\0&\textrm{ otherwise.}\end{cases}\]
Then we have
\[\sum_s g(s)\abs{A_1(s)}\abs{A_2(s)}
<\sum_s M\abs{A_1(s)}^{1/2}\abs{A_2(s)}^{1/2}.\]
To see why, note first that each summand on the left-hand side is $\leq$ the corresponding summand on the right-hand side, arguing by cases on whether $g(s)=1$ or not. It therefore suffices to show that there exists some $s$ such that the summand on the left-hand side is $<$ the corresponding summand on the right-hand side.

If $g(s)=0$ for all $s$ then choose some $s$ such that $\abs{A_1(s)}\abs{A_2(s)}\geq M^2$ (this must exist or else $\abs{A_1(s)}\abs{A_2(s)}=0$ for all $s$, but then $\norm{1_A\circ 1_A}_{p(\mu)}^p=0$ by the above calculation). Otherwise, choose some $s$ such that $g(s)=1$, and note that for this $s$ we have, by definition of $s$,
\[\abs{A_1(s)}\abs{A_2(s)}<M\abs{A_1(s)}^{1/2}\abs{A_2(s)}^{1/2}.\]

We now choose
\[M=\tfrac{1}{2}\abs{A}^{-p}(\abs{B_1}\abs{B_2})^{1/2}\norm{\ind{A}\circ \ind{A}}_{p(\mu)}^p,\]
so that, by the Cauchy-Schwarz inequality,
\begin{align*}
\sum_s g(s)\abs{A_1(s)}\abs{A_2(s)}
&< M \sum_s \abs{A_1(s)}^{1/2}\abs{A_2(s)}^{1/2}\\
&\leq M\brac{\sum_s \sum_{x\in G}1_{A_1(s)}(x)}^{1/2}\brac{\sum_s \sum_{x\in G}1_{A_2(s)}(x)}^{1/2}\\
&= M\abs{A}^p(\abs{B_1}\abs{B_2})^{1/2}\\
&=\frac{1}{2}\sum_s \abs{A_1(s)}\abs{A_2(s)}
\end{align*}
and so
\[
\sum_s (1-g(s))\abs{A_1(s)}\abs{A_2(s)}
> \frac{1}{2}\sum_s \abs{A_1(s)}\abs{A_2(s)}\]
whence
\[\sum_s\langle \ind{A_1(s)}\circ \ind{A_2(s)},f\rangle =\eta \sum \abs{A_1(s)}\abs{A_2(s)}
< 2\eta \sum_s \abs{A_1(s)}\abs{A_2(s)}(1-g(s)).\]
In particular there must exist some $s$ such that
\[\langle \ind{A_1(s)}\circ \ind{A_2(s)},f\rangle< 2\eta \abs{A_1(s)}\abs{A_2(s)}(1-g(s)).\]
We claim this $s$ meets the requirements. The first is satisfied since the right-hand side is $\leq 2\eta \abs{A_1(s)}\abs{A_2(s)}$. The second is satisfied since the left-hand side is trivially $\geq 0$ and hence such an $s$ must satisfy $g(s)=0$, whence either $ \abs{A_1(s)}\abs{A_2(s)}\geq M^2$, that is,
\[\abs{A_1(s)}\abs{A_2(s)}\geq \frac{1}{4}\abs{A}^{-2p}\abs{B_1}\abs{B_2}\norm{\ind{A}\circ \ind{A}}_{p(\mu)}^{2p},\]
or $\abs{A_1(s)}\abs{A_2(s)}=0$, which can't happen because then the right-hand side is $=0$.

The final bound now follows since $xy \leq \min(x,y)$ when $x,y\leq 1$.
\end{proof}


\begin{lemma}
\label{sift}
\lean{sifting}
\leanok
Let $\epsilon,\delta >0$ and $p\geq \max(2,\epsilon^{-1}\log(2/\delta))$ be an even integer. Let $B_1,B_2\subseteq G$, and let $\mu=\mu_{B_1}\circ\mu_{B_2}$. For any finite set $A\subseteq G$, if
\[S=\{ x\in G: \ind{A}\circ \ind{A}(x)>(1-\epsilon)\norm{\ind{A}\circ \ind{A}}_{p(\mu)}\},\]
then there are $A_1\subseteq B_1$ and $A_2\subseteq B_2$ such that
\[\langle \mu_{A_1}\circ\mu_{A_2},\ind{S}\rangle \geq 1-\delta\]
and
\[\min\brac{\frac{\abs{A_1}}{\abs{B_1}},\frac{\abs{A_2}}{\abs{B_2}}}\geq \frac{1}{4}\abs{A}^{-2p}\norm{\ind{A}\circ \ind{A}}_{p(\mu)}^{2p}.\]
\end{lemma}

\begin{proof}
\uses{drc}
\leanok
Apply Lemma~\ref{drc} with $f=\ind{G\backslash S}$. This produces some $A_1\subseteq B_1$ and $A_2\subseteq B_2$ such that
\[\langle \mu_{A_1}\circ \mu_{A_2}, \ind{G\backslash S}\rangle\leq 2\frac{\langle (\ind{A}\circ \ind{A})^p,\ind{G\backslash S}\rangle_\mu}{\norm{\ind{A}\circ \ind{A}}_{p(\mu)}^p}\]
and
\[\min\brac{\frac{\abs{A_1}}{\abs{B_1}},\frac{\abs{A_2}}{\abs{B_2}}}\geq \frac{1}{4}\abs{A}^{-2p}\norm{\ind{A}\circ \ind{A}}_{p(\mu)}^{2p}.\]
It then suffices to note that
\[\langle \mu_{A_1}\circ \mu_{A_2},\ind{S}\rangle=1-\langle \mu_{A_1}\circ \mu_{A_2},\ind{G\backslash S}\rangle\]
and by definition of $S$ we have
\[\langle (\ind{A}\circ \ind{A})^p,\ind{G\backslash S}\rangle_\mu\leq (1-\epsilon)^p\norm{\ind{A}\circ \ind{A}}_{p(\mu)}^p\sum_{x}\mu(x)=(1-\epsilon)^p\norm{\ind{A}\circ \ind{A}}_{p(\mu)}^p.\]
Now use the fact that $p\geq \epsilon^{-1}\log(2/\delta)$ together with the inequality $1-x\leq e^{-x}$ to deduce that the right-hand side is $\leq \tfrac{\delta}{2}\norm{\ind{A}\circ \ind{A}}_{p(\mu)}^p$.
\end{proof}


\begin{corollary}\label{sift_cor}
\lean{sifting_cor}
\leanok
Let $\epsilon,\delta>0$ and $p\geq \max(2,\epsilon^{-1}\log(2/\delta))$ be an even integer and $\mu\equiv 1/N$. If $A\subseteq G$ has density $\alpha$ and
\[S = \{ x : \mu_A\circ \mu_A(x) \geq (1-\epsilon)\| \mu_A\circ \mu_A\|_{p(\mu)}\}\]
then there are $A_1,A_2\subseteq G$ such that
\[\langle \mu_{A_1}\circ \mu_{A_2}, 1_S\rangle \geq 1-\delta\]
and both $A_1$ and $A_2$ have density
\[\geq \frac{1}{4}\alpha^{2p}.\]
\end{corollary}

\begin{proof}
\uses{sift}
\leanok
We apply Lemma~\ref{sift} with $B_1=B_2=G$. It remains to note that
\[\| 1_A\circ 1_A\|_{p(\mu)}\geq \| 1_A\circ 1_A\|_{1(\mu)}=\lvert A\rvert^2/N.\]
\end{proof}


%%FILE blueprint/src/chapter/ff.tex
\chapter{Finite field model}

\begin{theorem}
  \label{ap_in_ff}
\lean{ap_in_ff}
\leanok
If $A_1,A_2,S\subseteq \bbf_q^n$ are such that $A_1$ and $A_2$ both have density at least $\alpha$ then there is a subspace $V$ of codimension

\[\mathrm{codim}(V)\leq 2^{27}\lo(\alpha)^2\lo(\epsilon \alpha)^2\epsilon^{-2}\]
such that
\[\abs{\langle \mu_V\ast \mu_{A_1}\ast \mu_{A_2},\ind{S}\rangle-\langle \mu_{A_1}\ast \mu_{A_2},\ind{S}\rangle}\leq \epsilon.\]
\end{theorem}

\begin{proof}
\uses{chang, linfty_ap_boosted}
\leanok
(In this proof we write $G=\mathbb{F}_q^n$.) Let $k=\lceil \lo(\epsilon\alpha/4)\rceil$. Note that $\lvert A_1+G\rvert=\lvert G\rvert\leq \alpha^{-1}\lvert A\rvert$. Furthermore, $\lvert A_2\rvert/\lvert S\rvert\geq \alpha$. Therefore by Theorem~\ref{linfty_ap_boosted} there exists some $T\subseteq G$ with
\[\lvert T\rvert \geq \exp(-2^{16}\lo(\alpha)^2k^2\epsilon^{-2})\lvert S\rvert\]
such that
\[\| \mu_T^{(k)}\ast \mu_{A_1}\ast \mu_{A_2}\circ \ind{S}-\mu_{A_1}\ast \mu_{A_2}\circ \ind{S}\|_{\infty}\leq \epsilon/4.\]
Let $\Delta=\Delta_{1/2}(\mu_T)$ and
\[V = \{ x \in G : \gamma(x)=1\textrm{ for all }\gamma\in\Delta\}.\]
Note that
\[\langle \mu_V\ast \mu_{A_1}\ast \mu_{A_2},\ind{S}\rangle = \langle \mu_V,\mu_{A_1}\ast \mu_{A_2}\circ \ind{S}\rangle=\frac{1}{\abs{V}}\sum_{v\in V}\mu_{A_1}\ast \mu_{A_2}\circ \ind{S}(v)\]
and
\[\langle \mu_{A_1}\ast \mu_{A_2},\ind{S}\rangle = \mu_{A_1}\ast \mu_{A_2}\circ\ind{S}(0).\]
Therefore
\[\abs{\langle \mu_V\ast \mu_{A_1}\ast \mu_{A_2},\ind{S}\rangle-\langle \mu_{A_1}\ast \mu_{A_2},\ind{S}\rangle}\leq \frac{1}{\abs{V}}\sum_{v\in V}\abs{\mu_{A_1}\ast \mu_{A_2}\circ \ind{S}(v)-\mu_{A_1}\ast \mu_{A_2}\circ \ind{S}(0)}.\]
In particular it suffices to show that, for any $v\in V$,
\[\abs{\mu_{A_1}\ast \mu_{A_2}\circ \ind{S}(v)-\mu_{A_1}\ast \mu_{A_2}\circ \ind{S}(0)}\leq \epsilon.\]
By the triangle inequality and construction of $T$, it suffices to show that

\[\abs{\mu_T^{(k)}\ast \mu_{A_1}\ast \mu_{A_2}\circ \ind{S}(v)-\mu_T^{(k)}\ast\mu_{A_1}\ast \mu_{A_2}\circ \ind{S}(0)}\leq \epsilon/2.\]
By the Fourier transform we have, for any $x\in G$,
\[\mu_T^{(k)}\ast \mu_{A_1}\ast \mu_{A_2}\circ \ind{S}(x)=\frac{1}{N}\sum_{\gamma\in \widehat{G}}\widehat{\mu_T}(\gamma)^k\widehat{\mu_{A_1}}(\gamma)\widehat{\mu_{A_2}}(\gamma)\widehat{\ind{-S}}(\gamma)\gamma(x).\]
Therefore the left-hand side of the desired inequality is, by the triangle inequality, at most
\[\frac{1}{N}\sum_{\gamma\in \widehat{G}}\abs{\widehat{\mu_T}(\gamma)}^k\abs{\widehat{\mu_{A_1}}(\gamma)\widehat{\mu_{A_2}}(\gamma)\widehat{\ind{-S}}(\gamma)}\abs{\gamma(v)-1}.\]
By choice of $v\in V$ the summand vanishes when $\gamma\in\Delta$. When $\gamma\not\in \Delta$ the summand is bounded above by
\[2^{1-k}\abs{\widehat{\mu_{A_1}}(\gamma)\widehat{\mu_{A_2}}(\gamma)\widehat{\ind{-S}}(\gamma)}.\]
Therefore the left-hand side of the desired inequality is at most
\[2^{1-k}\frac{1}{N}\sum_{\gamma}\abs{\widehat{\mu_{A_1}}(\gamma)\widehat{\mu_{A_2}}(\gamma)\widehat{\ind{-S}}(\gamma)}\leq 2^{1-k}\abs{S}\frac{1}{N}\sum_{\gamma}\abs{\widehat{\mu_{A_1}}(\gamma)\widehat{\mu_{A_2}}(\gamma)}\]
using the trivial bound $\lvert\widehat{\ind{S}}\rvert\leq \abs{S}$. By the Cauchy-Schwarz inequality the sum on the right is at most
\[\brac{\sum_\gamma \abs{\widehat{\mu_{A_1}}}^2}^{1/2}\brac{\sum_\gamma \abs{\widehat{\mu_{A_2}}}^2}^{1/2}.\]
By Parseval's identity this is at most $\alpha^{-1}$. Therefore the desired inequality follows from
\[2^{1-k}\abs{S}\frac{1}{N}\alpha^{-1}\leq 2^{1-k}\alpha^{-1}\leq \epsilon/2.\]

It remains to check the codimension of $V$. For this, let $\Delta'\subseteq \Delta$ be as provided by Chang's lemma, Lemma~\ref{chang}, so that

\[\Delta\subseteq \left\{ \sum_{\gamma\in\Delta'}c_\gamma \gamma : c_\gamma\in \{-1,0,1\} \right\}.\]
Let

\[W= \{ x \in G : \gamma(x)=1\textrm{ for all }\gamma\in\Delta'\}.\]
It follows that $W\leq V$, so it suffices to bound the codimension of $W$. This we can bound trivially using the bound from Chang's lemma and the fact that $\lo(\delta)=\log(e^2/\delta)\leq 2+\log(1/\delta)\leq 4\log(1/\delta)$, provided $\log(1/\delta)\geq 1$, so
\[\abs{\Delta'}\leq  \lceil 4e\mathcal{L}(\delta)\rceil\leq 2^7\log(1/\delta),\]
where
\[\delta=\abs{T}/N\geq \exp(-2^{16}\lo(\alpha)^2k^2\epsilon^{-2}),\]
so
\[\mathrm{codim}(V)\leq \abs{\Delta'}\leq 2^{23}\lo(\alpha)^2k^2\epsilon^{-2}\leq 2^{25}\lo(\alpha)^2\lo(\epsilon \alpha/4)^2\epsilon^{-2},\]
and now use $\lo(\epsilon\alpha/4)\leq 2\lo(\epsilon\alpha)$, say.
\end{proof}


\begin{lemma}
\label{ast_le_circ}
\lean{cLpNorm_conv_le_cLpNorm_dconv}
\leanok
For any function $f:G\to \mathbb{C}$ and integer $k\geq 1$
\[\| f\ast f\|_{2k}\leq \| f\circ f\|_{2k}.\]
\end{lemma}

\begin{proof}
\leanok
To finish, similar trick to unbalancing.
\end{proof}


\begin{lemma}\label{ast_expand}
\lean{balance_conv}
\leanok
For any function $f$ with $\sum f(x)=1$
\[f\ast f-1/N = (f-1/N)\ast (f-1/N).\]
\end{lemma}

\begin{proof}
\leanok
Expand everything out.
\end{proof}


\begin{lemma}\label{circ_expand}
\lean{balance_dconv}
\leanok
For any function $f$ with $\sum f(x)=1$
\[f\circ f-1/N = (f-1/N)\circ (f-1/N).\]
\end{lemma}

\begin{proof}
\leanok
Expand everything out.
\end{proof}


\begin{lemma}
\label{global_dichotomy}
\lean{global_dichotomy}
\leanok
Let $\epsilon >0$ and $\mu\equiv 1/N$. If $A,C\subseteq G$, where $C$ has density at least $\gamma$, and
\[\abs{N\langle \mu_A\ast \mu_A,\mu_C\rangle -1}>\epsilon\]
then, if $f=(\mu_A-1/N)$, $\norm{f\circ f}_{p(\mu)} \geq \epsilon /2N$ for $p=2\lceil \lo(\gamma)\rceil$.
\end{lemma}

\begin{proof}
\uses{ast_le_circ, ast_expand, circ_expand}
\leanok
By H\"{o}lder's inequality, for any $p\geq 1$
\[\epsilon < \abs{ N\langle \mu_A \ast \mu_A - 1/N, \mu_C \rangle } \leq \norm{\mu_A\ast \mu_A-1/N}_p\gamma^{-1/p}N^{1-1/p}.\]
In particular if we choose $p=2\lceil\lo(\gamma)\rceil$ then $\gamma^{-1/p}\leq e^{1/2}\leq 2$ and so we deduce that, by Lemma~\ref{ast_expand},
\[\norm{f\ast f}_{p}\geq \tfrac{1}{2}\epsilon N^{1/p-1}.\]
It remains to use Lemmas~\ref{ast_expand} and \ref{circ_expand} and apply Lemma~\ref{ast_le_circ}, and note that we can pass from the $L^p$ norm to the $L^p(\mu)$ norm losing a factor of $N^{1/p}$.
\end{proof}


\begin{proposition}
\label{di_in_ff}
\lean{di_in_ff}
\leanok
Let $\epsilon \in (0,1)$. If $A,C\subseteq \bbf_q^n$, where $C$ has density at least $\gamma$, and
\[\abs{N\langle \mu_A\ast \mu_A,\mu_C\rangle -1}> \epsilon\]
then there is a subspace $V$ of codimension
\[\leq 2^{171}\epsilon^{-24}\lo(\alpha)^4\lo(\gamma)^4.\]
such that $\norm{\ind{A}\ast \mu_V}_\infty \geq (1+\epsilon/32)\alpha$.
\end{proposition}

\begin{proof}
\uses{global_dichotomy, unbalancing, ap_in_ff, circ_expand, sift_cor}
\leanok

By Lemma~\ref{global_dichotomy}, if $f=\mu_A-1/N$,
\[\norm{f\circ f}_{p(\mu)}\geq \epsilon /2N,\]
where $p=2\lceil \lo(\gamma)\rceil\leq 4\lo(\gamma)$. By Lemma~\ref{unbalancing} there exists some $p'$ such that
\[p'\leq 128\epsilon^{-1}\log(96/\epsilon)\lo(\gamma)\]
such that
\[\norm{f\circ f+1/N}_{p'(\mu)}\geq (1+\epsilon/4)/N.\]
By Lemma~\ref{circ_expand} $f\circ f+1/N=\mu_A\circ \mu_A$.

Let $q=2\lceil p'+2^8\epsilon^{-2}\log(64/\epsilon)\rceil$. By Corollary~\ref{sift_cor}, there are $A_1,A_2$, both of density $\geq \alpha^{2q}$ such that
\[\langle \mu_{A_1}\circ \mu_{A_2},\ind{S}\rangle \geq 1-\epsilon/32\]
where
\[S=\{x : \mu_A\circ \mu_A(x)\geq (1-\epsilon/16)\| \mu_A\circ \mu_A\|_{q(\mu)}\}.\]
Since
\[\| \mu_A\circ \mu_A\|_{q(\mu)}\geq \| \mu_A\circ \mu_A\|_{p'(\mu)}\geq (1+\epsilon/4)/N\]
we know
\[S\subseteq S'=\{x : \mu_A\circ \mu_A(x)\geq (1+\epsilon/8)/N\}.\]

By Theorem~\ref{ap_in_ff} (applied with $\epsilon$ replaced by $\epsilon/32$) there is a subspace $V$ of codimension
\[\leq 2^{37}\lo(\alpha^{2q})^2\lo(\epsilon\alpha^{2q}/32)^2\epsilon^{-2}\] such that
\[\langle \mu_V\ast \mu_{A_1}\circ \mu_{A_2},\ind{S'}\rangle \geq 1-\tfrac{1}{16}\epsilon.\]
Using $\lo(xy)\leq x^{-1}\lo(y)$ we have
\[\lo(\epsilon \alpha^{2q}/32)\leq 32\epsilon^{-1}\lo(\alpha^{2q}),\]
and we also use $\lo(x^y)\leq y\lo(x)$ to simplify the codimension bound to
\[\leq 2^{51}q^4\lo(\alpha)^4\epsilon^{-4}.\]
We further note that (using $\log x\leq x$ say)
\[q\leq 2^{10}p'\epsilon^{-2}\log(64/\epsilon)\leq 2^{30}\epsilon^{-5}\lo(\gamma).\]
Therefore the desired codimension bound follows. Finally, by definition of $S'$, it follows that
\begin{align*}
(1+\epsilon/32)/N
&\leq ((1+\epsilon/8)/N)(1-\epsilon/16)\\
&\leq \langle \mu_V\ast \mu_{A_1}\circ \mu_{A_2},\mu_A\circ \mu_A\rangle\\
&\leq \norm{\mu_V\ast \mu_A}_\infty \norm{\mu_{A}\ast \mu_{A_2}\circ \mu_{A_1}}_1\\
&= \norm{\mu_V\ast \ind{A}}_\infty\abs{A}^{-1},
\end{align*}
and the proof is complete.
\end{proof}


\begin{lemma}
  \label{no3aps_inner_prod}
  \lean{ThreeAPFree.wInner_one_mu_ddconv_mu_mu_two_smul_mu}
  \leanok

  If $A\subseteq G$ has no non-trivial three-term arithmetic progressions and $G$ has odd order then
  \[\langle \mu_A\ast \mu_A,\mu_{2\cdot A}\rangle = 1/\abs{A}^2.\]
\end{lemma}
\begin{proof}
  \leanok

  Expand out using definitions.
\end{proof}


\begin{theorem}
  \label{ff}
  \uses{}
  \lean{ff}
  \leanok

  Let $q$ be an odd prime power. If $A\subseteq \mathbb{F}_q^n$ with $\alpha=\abs{A}/q^n$ has no non-trivial three-term arithmetic progressions then
  \[n \ll \lo(\alpha)^9.\]
\end{theorem}
\begin{proof}
  \uses{no3aps_inner_prod, di_in_ff}
  \leanok

  Let $t\geq 0$ be maximal such that there is a sequence of subspaces $\bbf_q^n=V_0\geq \cdots \geq V_t$ and associated $A_i\subseteq V_i$ with $A_0=A$ such that
  \begin{enumerate}
  \item for $0\leq i\leq t$ there exists $x_i$ such that $A_i\subseteq A-x_i$,
  \item with $\alpha_i=\abs{A_i}/\abs{V_i}$ we have
  \[\alpha_{i+1}\geq \frac{65}{64}\alpha_i\]
  for $0\leq i<t$, and
  \item \[\mathrm{codim} (V_{i+1}) \leq \mathrm{codim}(V_i)+O(\lo(\alpha)^8)\]
  for $0\leq i<t$. (here the second summand should be replaced with whatever explicit codimension bound we get from the above).
  \end{enumerate}

  Note this is well-defined since $t=0$ meets the requirements, and this process is finite and $t \ll \lo(\alpha)$, since $\alpha_i\leq 1$ for all $i$. Therefore
  \[\mathrm{codim}(V_t) \ll \lo(\alpha)^9.\]

  Suppose first that
  \[\lvert V_t\rvert\langle \mu_{A_t}\ast \mu_{A_t},\mu_{2\cdot A_t}\rangle<1/2.\]
  In this case we now apply Proposition~\ref{di_in_ff} to $A_t\subseteq V_t$ with $\epsilon=1/2$ (note that $N=\lvert V_t\rvert$ and all inner product, $\mu$ etc, are relative to the ambient group $V_t$ now). Therefore there is a subspace $V\leq V_t$ of codimension (relative to $V_t$) $\ll \lo(\alpha)^8$ such that there exists some $x\in V_t$ with
  \[\frac{\lvert (A_t-x)\cap V\rvert}{\lvert V\rvert}=1_{A_t}\ast \mu_V(x)=\| 1_{A_t}\ast \mu_V\|_\infty \geq (1+1/64)\alpha_t,\]
  which contradicts the maximality of $t$, letting $V_{t+1}=V$ and $A_{t+1}=(A_t-x)\cap V_t$.

  Therefore

  \[\lvert V_t\rvert\langle \mu_{A_t}\ast \mu_{A_t},\mu_{2\cdot A_t}\rangle\geq 1/2.\]
  By Lemma~\ref{no3aps_inner_prod} the left-hand side is equal to $\lvert V_t\rvert/\lvert A_t\rvert^2$, and therefore
  \[\alpha^2 \leq \alpha_t^2 \leq 2/\lvert V_t\rvert.\]
  By the codimension bound the right-hand side is at most
  \[2q^{O(\lo(\alpha)^9)-n}.\]
  If $\alpha^2 \leq 2q^{-n/2}$ we are done, otherwise we deduce that $\lo(\alpha)^9 \gg n$ as required.
\end{proof}


%%FILE blueprint/src/chapter/integers.tex
\chapter{The integer case}

\begin{theorem}
\label{th-ap-int}
\uses{bohr-reg-def}
There is a constant $c>0$ such that the following holds. Let $\epsilon>0$ and $B,B'\subseteq G$ be regular Bohr sets of rank $d$. Suppose that $A_1\subseteq B$ with density $\alpha_1$ and $A_2$ is such that there exists $x$ with $A_2\subseteq B'-x$ with density $\alpha_2$. Let $S$ be any set with $\abs{S}\leq 2\abs{B}$. There is a regular Bohr set $B''\subseteq B'$ of rank at most
\[d+O_\epsilon(\lo{\alpha_1}^3\lo{\alpha_2})\]
and size
\[\abs{B''}\geq \exp(-O_\epsilon(d\lo{\alpha_1\alpha_2/d}+\lo{\alpha_1}^3\lo{\alpha_2}\lo{\alpha_1\alpha_2/d}))\abs{B'}\]
such that
\[\abs{\langle \mu_{B'}\ast \mu_{A_1}\circ \mu_{A_2},\ind{S}\rangle-\langle \mu_{A_1}\circ \mu_{A_2},\ind{S}\rangle}\leq \epsilon.\]
\end{theorem}
\begin{proof}
\uses{linfty_ap, chang, bohr-size}
To do.
\end{proof}

\begin{proposition}
\label{Lp-orth}
\uses{bohr-reg-def}
There is a constant $c>0$ such that the following holds. Let $\epsilon >0$ and $p \geq 2$ be an integer. Let $B \subseteq G$ be a regular Bohr set and $A\subseteq B$ with relative density $\alpha$. Let $\nu : G \to \bbr_{\geq 0}$ be supported on $B_\rho$, where $\rho \leq c\epsilon\alpha/\rk(B)$, such that $\norm{\nu}_1=1$ and $\widehat{\nu}\geq 0$. If
    \[ \norm{(\mu_A-\mu_B) \circ (\mu_{A}-\mu_B)}_{p(\nu)} \geq \epsilon\, \mu(B)^{-1}, \]
    then there exists $p'\ll_\epsilon p$ such that
    \[ \norm{ \mu_{A}\circ \mu_{A}}_{p'(\nu)} \geq \left(1+\tfrac{1}{4}\epsilon\right) \mu(B)^{-1}. \]
\end{proposition}
\begin{proof}
\uses{reg-conv, unbalancing}
To do.
\end{proof}


\begin{proposition}
\label{pos-def-measures}
\uses{bohr-reg-def}
There is a constant $c>0$ such that the following holds. Let $p \geq 2$ be an even integer. Let $f : G \to \bbr$, let $B \subseteq G$ and $B', B'' \subseteq B_{c/\rk(B)}$ all be regular Bohr sets. Then
\[ \norm{ f\circ f }_{p(\mu_{B'}\circ \mu_{B'}\ast \mu_{B''}\circ \mu_{B''})} \geq \tfrac{1}{2} \norm{f*f}_{p(\mu_B)}. \]
\end{proposition}
\begin{proof}
\uses{bohr-majorise}
To do,
\end{proof}

\begin{proposition}
\label{hoelder-lifting}
\uses{bohr-reg-def}
There is a constant $c>0$ such that the following holds. Let $\epsilon >0$. Let $B \subseteq G$ be a regular Bohr set and $A\subseteq B$ with relative density $\alpha$, and let $B' \subseteq B_{c\epsilon\alpha/\rk(B)}$ be a regular Bohr set and $C\subseteq B'$ with relative density $\gamma$. Either
\begin{enumerate}
\item $\abs{ \langle \mu_A*\mu_A, \mu_{C} \rangle - \mu(B)^{-1} } \leq \epsilon \mu(B)^{-1}$ or
\item there is some $p \ll\lo{\gamma}$ such that $\norm{ (\mu_A-\mu_B)*(\mu_A-\mu_B)}_{p(\mu_{B'})} \geq \tfrac{1}{2}\epsilon \mu(B)^{-1}$.
\end{enumerate}
\end{proposition}
\begin{proof}
%\uses{reg-conv}
To do.
\end{proof}

\begin{proposition}
\label{prop-it}
\uses{bohr-reg-def}
There is a constant $c>0$ such that the following holds. Let $\epsilon>0$ and $p,k\geq 1$ be integers such that $(k,\abs{G})=1$. Let $B,B',B''\subseteq G$ be regular Bohr sets of rank $d$ such that $B''\subseteq B'_{c/d}$ and $A\subseteq B$ with relative density $\alpha$. If
    \[ \norm{ \mu_{A}\circ \mu_{A}}_{p(\mu_{k\cdot B'}\circ\mu_{k\cdot B'}\ast \mu_{k\cdot B''}\circ \mu_{k\cdot B''})} \geq \left(1+\epsilon\right) \mu(B)^{-1},\]
    then there is a regular Bohr set $B'''\subseteq B''$ of rank at most
    \[\rk(B''')\leq d+O_{\epsilon}(\lo{\alpha}^4p^4)\]
    and size
    \[\abs{B'''}\geq \exp(-O_{\epsilon}(dp\lo{\alpha/d}+\lo{\alpha}^5p^5))\abs{B''}\]
    such that
    \[ \norm{ \mu_{B'''}*\mu_A }_\infty \geq (1+c\epsilon)\mu(B)^{-1}. \]
\end{proposition}
\begin{proof}
\uses{sift, th-ap-int}
To do.
\end{proof}

\begin{theorem}
\label{flat-on-bohr}
\uses{bohr-reg-def}
There is a constant $c>0$ such that the following holds. Let ${\epsilon,\delta\in (0,1)}$ and $p,k\geq 1$ be integers such that $(k,\abs{G})=1$. For any $A\subseteq G$ with density $\alpha$ there is a regular Bohr set $B$ with
\[ d=\rk(B) =O_{\epsilon}\left(\lo{\alpha}^5p^4\right) \quad\text{and}\quad \abs{B}\geq \exp\left(-O_{\epsilon,\delta}(\lo{\alpha}^6p^5\lo{\alpha/p})\right)\abs{G} \]
and some $A'\subseteq (A-x)\cap B$ for some $x \in G$ such that
\begin{enumerate}
\item $\abs{A'}\geq (1-\epsilon)\alpha\abs{B}$,
\item $\abs{A'\cap B'}\geq (1-\epsilon)\alpha\abs{B'}$, where $B'=B_{\rho}$ is a regular Bohr set with ${\rho\in (\tfrac{1}{2},1)\cdot c\delta\alpha/d}$, and
\item
\[\norm{\mu_{A'}\circ \mu_{A'}}_{p(\mu_{k\cdot B''}\circ \mu_{k\cdot B''}\ast \mu_{k\cdot B'''}\circ \mu_{k\cdot B'''})} <(1+ \epsilon)\mu(B)^{-1},\]
for any regular Bohr sets $B'' = B'_{\rho'}$ and $B'''=B''_{\rho''}$ satisfying ${\rho',\rho''\in(\frac{1}{2},1)\cdot c\delta\alpha/d}$.
\end{enumerate}
\end{theorem}
\begin{proof}
\uses{prop-it, bourgain-trick, bohr-size}
To do.
\end{proof}

\begin{theorem}
\label{th-int-gen}
\uses{bohr-reg-def}
There is a constant $c>0$ such that the following holds. Let $\delta,\epsilon\in (0,1)$, let $p \geq 1$ and let $k$ be a positive integer such that $(k,\abs{G})=1$. There is a constant $C=C(\epsilon,\delta,k)>0$ such that the following holds.

For any finite abelian group $G$ and any subset $A\subseteq G$ with $\abs{A}=\alpha \abs{G}$ there exists a regular Bohr set $B$ with
\[\rk(B)\leq Cp^4\log(2/\alpha)^5\]
and
\[\abs{B}\geq \exp\left(-Cp^5\log(2p/\alpha)\log(2/\alpha)^6\right)\abs{G}\]
and $A' \subseteq (A-x)\cap B$ for some $x\in G$ such that
\begin{enumerate}
\item $\abs{A'}\geq (1-\epsilon)\alpha \abs{B}$,
\item $\abs{A'\cap B'}\geq (1-\epsilon)\alpha\abs{B'}$, where $B'=B_{\rho}$ is a regular Bohr set with $\rho\in (\tfrac{1}{2},1)\cdot c\delta\alpha/dk$, and
\item
\[\norm{(\mu_{A'}-\mu_B)\ast (\mu_{A'}- \mu_B)}_{p(\mu_{k\cdot B'})} \leq \epsilon\frac{\abs{G}}{\abs{B}}.\]
\end{enumerate}
\end{theorem}
\begin{proof}
\uses{Lp-orth, flat-on-bohr, pos-def-measures}
To do.
\end{proof}

\begin{theorem}\label{main-int-count}
If $A\subseteq \{1,\ldots,N\}$ has size $\abs{A}=\alpha N$, then $A$ contains at least
\[\exp(-O(\lo{\alpha}^{12}))N^2\]
many three-term arithmetic progressions.
\end{theorem}
\begin{proof}
\uses{th-int-gen, hoelder-lifting}
To do.
\end{proof}

\begin{theorem}[Integer case]
  \label{int}
  \uses{}
  \lean{int}
  \leanok
  If $A\subseteq \{1,\ldots,N\}$ contains no non-trivial three-term arithmetic progressions then
  \[\lvert A\rvert \leq \frac{N}{\exp(-c(\log N)^{1/12})}\]
\end{theorem}
\begin{proof}
\uses{main-int-count}
  To do.
\end{proof}


%%FILE blueprint/src/chapter/main.tex
% This is the main point of entry to the blueprint.
% Add chapters of the blueprint here.
% This file is not meant to be built. Build src/web.tex or src/print.tex instead.

\input{chapter/ap}
\input{chapter/chang}
\input{chapter/unbalancing}
\input{chapter/drc}
\input{chapter/ff}
\input{chapter/bohr}
\input{chapter/integers}


%%FILE blueprint/src/chapter/unbalancing.tex
\chapter{Unbalancing}

\begin{lemma}
\label{nonneg_moments}
\lean{pow_inner_nonneg}
\leanok
For any function $f:G\to\bbr$ and integer $k\geq 0$
\[\bbe_x f\circ f(x)^k\geq 0.\]
\end{lemma}
\begin{proof}
\leanok
Test.
\end{proof}

\begin{lemma}
\label{unbalancing}
\lean{unbalancing}
\leanok
Let $\epsilon\in(0,1)$ and $\nu:G\to\bbr_{\geq 0}$ be some probability measure such that $\widehat{\nu}\geq 0$. Let $f:G\to\bbr$. If $\norm{f\circ f}_{p(\nu)}\geq \epsilon$ for some $p\geq 1$ then $\norm{f\circ f+1}_{p'(\nu)}\geq 1+\tfrac{1}{2}\epsilon$ for $p'=120\epsilon^{-1}\log(3/\epsilon)$.
\end{lemma}

\begin{proof}
\uses{nonneg_moments}
\leanok
Up to gaining a factor of 5 in $p'$, we can assume that $p\geq 5$ is an odd integer. % Want (3/4)^p \leq 1/4
Since the Fourier transforms of $f$ and $\nu$ are non-negative,
\[\bbe \nu f^{p}= \widehat{\nu}\ast\widehat{f}^{(p)}(0_{\widehat{G}})\geq 0.\]
It follows that, since $2\max(x,0)=x+\abs{x}$ for $x\in\bbr$,
\[2\langle \max(f,0),f^{p-1}\rangle_\nu=\bbe \nu f^{p}+\langle \abs{f},f^{p-1}\rangle_\nu\geq \norm{f}_{p(\nu)}^p\geq \epsilon^p.\]
Therefore, if $P=\{ x : f(x) \geq 0\}$, then $\langle \ind{P},f^{p}\rangle_\nu\geq \frac{1}{2}\epsilon^{p}$. Furthermore, if $T=\{ x\in P : f(x) \geq \tfrac{3}{4}\epsilon\}$ then $\langle \ind{P\backslash T},f^p\rangle_\nu \leq \tfrac{1}{4}\epsilon^{p}$, and hence by the Cauchy-Schwarz inequality,
\[\nu(T)^{1/2}\norm{f}_{2p(\nu)}^{p}\geq \langle \ind{T}, f^{p}\rangle_\nu \geq \tfrac{1}{4}\epsilon^{p}.\]
On the other hand, by the triangle inequality
\[\norm{f}_{2p(\nu)}\leq 1+\norm{f+1}_{2p(\nu)}\leq 3,\]
or else we are done, with $p'=2p$. Hence $\nu(T)\geq (\epsilon/3)^{3p}$. It follows that, for any $p'\geq 1$,
\[\norm{f+1}_{p'(\nu)}\geq \langle \ind{T},\abs{f+1}^{p'}\rangle_\nu^{1/p'}\geq  (1+\tfrac{3}{4}\epsilon)(\epsilon/3)^{3p/p'}.\]
The desired bound now follows if we choose $p'=24\epsilon^{-1}\log(3/\epsilon)p$, using $1-x\leq e^{-x}$.
\end{proof}

